% Vertical version mark in the right margin, at footer height.
% For the experimenter; easy to miss for participants.
\newcommand{\hiddenversion}[1]{%
  \par\noindent
  \makebox[0pt][l]{%
    \hspace*{\dimexpr\paperwidth-1in-\hoffset
      -\ifodd\value{page}\oddsidemargin\else\evensidemargin\fi
    -8mm\relax}%
    \raisebox{-\dimexpr\pagegoal-\pagetotal+22mm\relax}[0pt][0pt]{%
      \rotatebox{90}{%
        \fontsize{6}{7}\selectfont\textcolor{black!45}{#1}%
      }%
    }%
  }%
  \par
}

\section{Hidden Profile Exercise}
\label{annex:hidden-profile-exercise}

Diese Dokumente basieren auf der von Baker \cite{bakerEnhancingGroupDecision2010} vorgeschlagenen
Hidden-Profile-Aufgabe (Personalauswahl unter drei Kandidaten). Für die vorliegende Studie wurden die
Materialien ins Deutsche übersetzt. Geografische Angaben wurden angepasst.

Die nachfolgenden Seiten sind die Handouts, die den Teilnehmenden in der Sitzung direkt ausgehändigt
wurden. Damit die Versuchsleitung die unterschiedlich verteilten Informationsstände zuordnen kann,
ohne dass die Teilnehmenden merken, dass sie unterschiedliche Informationen erhalten, trägt jede Seite
am rechten unteren Rand eine unscheinbare, vertikal gesetzte Versionskennung
(z.\,B.\ \texttt{V.2026K01-01/03}).
Das Präfix \texttt{V.2026} ist absichtlich bedeutungslos und dient nur der Tarnung, damit die Kennung
nicht wie eine systematische Versuchskodierung wirkt.
\texttt{K0X} gibt an, an welchen der vier Teilnehmenden das Blatt ausgehändigt wird (\texttt{K01} bis
\texttt{K04}).
\texttt{0X/03} bezeichnet die Informationen zum \texttt{X}-ten von insgesamt drei Kandidaten, zwischen
denen die Gruppe wählen soll.
Eine offene, blattspezifische Nummer wäre verdächtig gewesen und hätte die für Hidden Profiles
notwendige Informationsasymmetrie gefährdet. Die Markierung ist daher bewusst klein, hell und
außerhalb des Lesebereichs platziert; sie dient ausschließlich der internen Zuordnung.

Aus demselben Grund wurden auf diesen Seiten auch die Seitenzahlen im Footer weggelassen. (Zur vermeindung
  dass während der Interviews es zu absprachen kommt wie: "Auf Seite X steht ja, dass...?" Dies könnte bei
unterschiedlichen Seitenzahlen den teilnehmden verraten, dass sie unterschiedliche Blätter erhalten haben.)

\clearpage

% ---------------------------------------------------------------------------------------------------------------------
\subsection*{Aufgabenstellung und Stellenausschreibung}
\pagestyle{empty}
\thispagestyle{empty}

\textbf{Aufgabenstellung:}

Sie sind Mitglied des Auswahlgremiums für einen neuen Präsidenten der Higher Education University, einer
privaten Liberal-Arts-Hochschule in Georgia. Nachstehend finden Sie die wichtigsten Punkte aus der
Stellenausschreibung. Es haben sich drei Kandidaten beworben; die Informationen, die Sie aus den
Lebensläufen, Referenzprüfungen, Vorstellungsgesprächen und Präsentationen der Kandidaten gewonnen haben,
sind auf den folgenden Seiten aufgeführt. Lesen Sie das gesamte Material durch und bestimmen Sie anschließend
den Kandidaten, dem Ihrer Meinung nach die Stelle angeboten werden sollte. Stellenausschreibung für den
Präsidenten der Higher Education University.
Die Higher Education University ist eine kleine (2.500 Studierende) private Liberal-Arts-Hochschule in
Georgia.
Der Präsident der Universität trägt die letztendliche Verantwortung für die Verwaltung und den
Betrieb der Universität. Er beaufsichtigt die strategische Planung, das Fundraising und die Haushaltsplanung,
stellt sicher, dass die Richtlinien der Universität eingehalten werden, und pflegt die Beziehungen zur
Gemeinde, zu den Alumni und zu anderen externen Interessengruppen.

\clearpage

\textbf{Gewünsche Qualitfikationen:}

Die Bewerber werden anhand der folgenden Kriterien bewertet:
\begin{itemize}
  \item Nachweis akademischer Leistungen und Lehrerfahrung oder einschlägiger, umfassender beruflicher
    Qualifikationen, die die Fähigkeit zur Führung in einem akademischen Umfeld belegen;
  \item Nachweis von Führungsqualitäten in einem akademischen oder vergleichbaren Umfeld mit den folgenden Merkmalen:
    \begin{itemize}
      \item Integrität und Bekenntnis zu höchsten ethischen Standards;
      \item Engagement für eine partizipative Leitung und die unfassende Einbindung von Lehrenden,
        Mitarbeitern und Studierenden;
      \item Engagement für eine vielfältige Gemeinschaft aus Lehrenden, Mitarbeitern und Studierenden mit
        Chancengleichheit für alle;
    \end{itemize}
  \item Die Fähigkeit, eine Vision für die Universität zu entwickeln und die Führung und Inspiration zu
    bieten, die erforderlich sind, um diese Vision zu verwirklichen;
  \item Das Engagement für die Schaffung einer Kultur, die herausragende Leistungen in Forschung, Lehre und
    Service wertschätzt;
  \item Die Fähigkeit, die Universität gegenüber externen Interessengruppen zu präsentieren und die
    Unterstützung für ihren Auftrag und ihre Programme zu stärken;
  \item Das Engagement für die Verbesserung der Lebensqualität von Studierenden, Lehrenden und Mitarbeitern;
  \item Die Unterstützung bei der Entwicklung und Nutzung neuer Technologien zur Förderung des Lernens, der
    wissenschaftlichen Arbeit und des Ideenaustauschs.
\end{itemize}

\clearpage
% ---------------------------------------------------------------------------------------------------------------------
\subsection*{Kandidateninformationen - Kandidat Stevens}
\hiddenversion{V.2026K01-01/03}

\clearpage
\pagestyle{TUDa}
\thispagestyle{TUDa}
% ---------------------------------------------------------------------------------------------------------------------
\subsection*{Qualifikationen und Eigenschaften von Kandidaten für das Amt des Präsidenten einer Hochschule}